\begin{theorem}[Frame correspondence and completeness]
In each row, let \(\mathsf L\) be the displayed calculus and let
\(\mathcal C_{\mathsf L}\) be the displayed class of frames:
\[
\begin{array}{c|c}
  \text{calculus}&\text{frame class}\\ \hline
  \K&\text{all frames}\\
  \mathsf T&\text{reflexive frames}\\
  \Sfour&\text{reflexive transitive frames}\\
  \Sfive&\text{frames whose accessibility relation is an equivalence relation.}
\end{array}
\]
Then, for every premise set \(\Gamma\) and formula \(A\),
\[
  \Gamma\derives_{\mathsf L}A
  \quad\Longleftrightarrow\quad
  \Gamma\entails^{\ell}_{\mathcal C_{\mathsf L}}A.
\]
Thus each calculus is sound and strongly complete for local consequence on its
paired frame class.
\end{theorem}

\paragraph{Proof sketch.}
Soundness checks each axiom against its frame condition and verifies modus
ponens and necessitation.  Completeness constructs the canonical structure of
maximally consistent sets.  Define \(wRv\) when every \(\Box A\in w\) has
\(A\in v\), and prove the truth lemma by induction.  The axioms T, 4, and 5
force the canonical relation to be reflexive, transitive, and Euclidean;
reflexivity plus Euclideanness yields an equivalence relation.  See
\cite{BlackburnDeRijkeVenema2001} for the detailed canonical-structure
arguments.

